\appendix
\appendixpage

\chapter{Code Listing}
\label{chap:code}

\begin{minted}
[
frame=lines,
framesep=2mm,
baselinestretch=1.0,
fontsize=\footnotesize,
linenos
]
{python}
#!/usr/bin/env python3
# bernstein_polynomial.py
import math
import sys
import numpy as np

def bernstein_polynomial(i, p, nti):
    """ 
    Computes the Bernstein polynomial coefficient for 
    control point i with 
    polynomial degreee p
    for a 1D parameter array t in interval [0, 1] broken into 
    nti number of equidistance time intervals.
    """
    if i >= 0 and p >= 1 and nti >= 2 and i <= p:
        t = np.linspace(0, 1, nti+1)
        bp = math.factorial(p) / \
            (math.factorial(i) * math.factorial(p-i)) * t**i * (1-t)**(p-i)
        return bp
    else:
        print(f'Input (i, p, nti) = ({i}, {p}, {nti}) is out of range.')
        print('i is non-negative integer 0, 1, 2, ... p')
        print('p is integer >= 1')
        print('nti is integer >= 2')
        return None
\end{minted}

\begin{minted}
[
frame=lines,
framesep=2mm,
baselinestretch=1.0,
fontsize=\footnotesize,
linenos
]
{python}
#!/usr/bin/env python3
# bernstein_polynomial_test.py
"""
This module is a unit test of the bernstein_polynomial implementation.

To run
$ python bernstein_polynomial_test.py              # terse interaction
$ python -m unittest bernstein_polynomial_test     # default interaction
$ python -m unittest -v bernstein_polynomial_test  # verbose interaction
"""
# standard library imports
import sys
# 
import numpy as np
from unittest import TestCase, main
# 
import bernstein_polynomial as bp

class TestBernstein(TestCase):

    @classmethod
    def setUpClass(cls):
        cls._TOL = 1e-6  # tolerance
        cls._nti = 4  # number of time intervals
        # t, e.g., four intervals, five evaluation points
        cls._t = np.linspace(0, 1, cls._nti+1)
        cls._verbosity = 0

    @classmethod
    def same(cls, a, b):
        same_to_tolerance = False
        l2norm_diff = np.linalg.norm(a - b)

        if cls._verbosity > 0:
            print(f'array a = {a}')
            print(f'array b = {b}')
            print(f'l2norm_diff = {l2norm_diff}')

        if np.abs(l2norm_diff) < cls._TOL:
            same_to_tolerance = True

        return same_to_tolerance

    def test_000_b01(self):
        known = 1 - self._t
        i, p = 0, 1
        calc = bp.bernstein_polynomial(i, p, self._nti)
        self.assertTrue(self.same(known, calc))

    def test_001_b11(self):
        known = self._t
        i, p = 1, 1
        calc = bp.bernstein_polynomial(i, p, self._nti)
        self.assertTrue(self.same(known, calc))

    def test_003_b02(self):
        known = (1 - self._t)**2
        i, p = 0, 2
        calc = bp.bernstein_polynomial(i, p, self._nti)
        self.assertTrue(self.same(known, calc))

    def test_004_b12(self):
        known = 2 * self._t * (1 - self._t)
        i, p = 1, 2
        calc = bp.bernstein_polynomial(i, p, self._nti)
        self.assertTrue(self.same(known, calc))
        
    def test_005_b22(self):
        known = (self._t)**2
        i, p = 2, 2
        calc = bp.bernstein_polynomial(i, p, self._nti)
        self.assertTrue(self.same(known, calc))

    def test_006_b03(self):
        known = (1 - self._t)**3
        i, p = 0, 3
        calc = bp.bernstein_polynomial(i, p, self._nti)
        self.assertTrue(self.same(known, calc))

    def test_007_b13(self):
        known = 3 * self._t * (1 - self._t)**2
        i, p = 1, 3
        calc = bp.bernstein_polynomial(i, p, self._nti)
        self.assertTrue(self.same(known, calc))

    def test_008_b23(self):
        known = 3 * (self._t)**2 * (1 - self._t)
        i, p = 2, 3
        calc = bp.bernstein_polynomial(i, p, self._nti)
        self.assertTrue(self.same(known, calc))
        
    def test_009_b33(self):
        known = (self._t)**3
        i, p = 3, 3
        calc = bp.bernstein_polynomial(i, p, self._nti)
        self.assertTrue(self.same(known, calc))

    def test_010_b04(self):
        known = (1 - self._t)**4
        i, p = 0, 4
        calc = bp.bernstein_polynomial(i, p, self._nti)
        self.assertTrue(self.same(known, calc))

    def test_011_b14(self):
        known = 4 * self._t * (1 - self._t)**3
        i, p = 1, 4
        calc = bp.bernstein_polynomial(i, p, self._nti)
        self.assertTrue(self.same(known, calc))

    def test_012_b24(self):
        known = 6 * (self._t)**2 * (1 - self._t)**2
        i, p = 2, 4
        calc = bp.bernstein_polynomial(i, p, self._nti)
        self.assertTrue(self.same(known, calc))

    def test_013_b34(self):
        known = 4 * (self._t)**3 * (1 - self._t)
        i, p = 3, 4
        calc = bp.bernstein_polynomial(i, p, self._nti)
        self.assertTrue(self.same(known, calc))

    def test_014_b44(self):
        known = (self._t)**4
        i, p = 4, 4
        calc = bp.bernstein_polynomial(i, p, self._nti)
        self.assertTrue(self.same(known, calc))

if __name__ == '__main__':
    main()  # calls unittest.main()
\end{minted}

\begin{minted}
[
frame=lines,
framesep=2mm,
baselinestretch=1.0,
fontsize=\footnotesize,
linenos
]
{python}
#!/usr/bin/env python3
# bernstein_extended.py
import numpy as np
import matplotlib.pyplot as plt
from matplotlib import rc
from matplotlib.ticker import AutoMinorLocator, MultipleLocator

import bernstein_polynomial as bp

LATEX = 1
SERIALIZE = 0

if LATEX:
  rc('font', **{'family': 'serif', 'serif': ['Computer Modern Roman']})
  rc('text', usetex=True)

nts = 101 # number of time steps, include t0 as a step
nti = nts - 1 # number of time intervals
t = np.linspace(0, 1, nts)  # parameterization

# fig = plt.figure(figsize=(6, 6))  # x in inches wide, y inches tall
fig = plt.figure(figsize=(6.5, 6.5))  # x in incmes wide, y inches tall
ax2 = fig.add_subplot(2, 2, 3, aspect=1)
ax0 = fig.add_subplot(2, 2, 1, aspect=1)
ax3 = fig.add_subplot(2, 2, 4, aspect=1, sharey=ax2)
ax1 = fig.add_subplot(2, 2, 2, aspect=1, sharey=ax0, sharex=ax3)

# globals
ap = dict(arrowstyle='->')  # arrow properties
lax, lay = 0.5, 1.1
b0_x = 0.32
b0_y = 0.87
dx = 0.008
annotate_dict = dict(arrowprops=ap, 
    ha='center', 
    va='bottom',
    backgroundcolor='white')
text_dict = dict(ha='center', 
    va='bottom',
    backgroundcolor='white')
linestyle_tuple = ('solid', 'dashed', 'dashdot')

p = 5
for q in np.arange(p+1):
    ls = linestyle_tuple[q % len(linestyle_tuple)]
    ax0.plot(t, bp.bernstein_polynomial(q, p, nti), linestyle=ls)

ax0.annotate(r'$b_{0,5}(t)$', xy=(0.06-dx, 0.75), xytext=(b0_x, b0_y), **annotate_dict)
ax0.text(lax, b0_y, r'$\ldots$', **text_dict)
ax0.annotate(r'$b_{5,5}(t)$', xy=(0.94+dx, 0.75), xytext=(1-b0_x, b0_y), **annotate_dict)

ax0.xaxis.set_major_locator(MultipleLocator(0.25))
ax0.yaxis.set_major_locator(MultipleLocator(0.25))
ax0.set_ylabel(r'$b_{i,p}(t)$')
ax0.text(lax, lay, r'(a) $p=5$', **text_dict)
ax0.grid()
plt.setp(ax0.get_xticklabels(), visible=False)

p = 6
for q in np.arange(p+1):
    ls = linestyle_tuple[q % len(linestyle_tuple)]
    ax1.plot(t, bp.bernstein_polynomial(q, p, nti), linestyle=ls)

ax1.annotate(r'$b_{0,6}(t)$', xy=(0.06-2*dx, 0.75), xytext=(b0_x, b0_y), **annotate_dict)
ax1.text(lax, b0_y, r'$\ldots$', **text_dict)
ax1.annotate(r'$b_{6,6}(t)$', xy=(0.94+2*dx, 0.75), xytext=(1-b0_x, b0_y), **annotate_dict)
ax1.xaxis.set_major_locator(MultipleLocator(0.25))
ax1.yaxis.set_major_locator(MultipleLocator(0.25))
ax1.text(lax, lay, r'(b) $p=6$', **text_dict)
ax1.grid()
plt.setp(ax1.get_xticklabels(), visible=False)
plt.setp(ax1.get_yticklabels(), visible=False)

p = 7
for q in np.arange(p+1):
    ls = linestyle_tuple[q % len(linestyle_tuple)]
    ax2.plot(t, bp.bernstein_polynomial(q, p, nti), linestyle=ls)

ax2.annotate(r'$b_{0,7}(t)$', xy=(0.06-3*dx, 0.75), xytext=(b0_x, b0_y), **annotate_dict)
ax2.text(lax, b0_y, r'$\ldots$', **text_dict)
ax2.annotate(r'$b_{7,7}(t)$', xy=(0.94+3*dx, 0.75), xytext=(1-b0_x, b0_y), **annotate_dict)

ax2.xaxis.set_major_locator(MultipleLocator(0.25))
ax2.yaxis.set_major_locator(MultipleLocator(0.25))
ax2.set_xlabel(r'$t$')
ax2.set_ylabel(r'$b_{i,p}(t)$')
ax2.text(lax, lay, r'(c) $p=7$', **text_dict)
ax2.grid()

p = 8
for q in np.arange(p+1):
    ls = linestyle_tuple[q % len(linestyle_tuple)]
    ax3.plot(t, bp.bernstein_polynomial(q, p, nti), linestyle=ls)

ax3.annotate(r'$b_{0,8}(t)$', xy=(0.06-4*dx, 0.75), xytext=(b0_x, b0_y), **annotate_dict)
ax3.text(lax, b0_y, r'$\ldots$', **text_dict)
ax3.annotate(r'$b_{8,8}(t)$', xy=(0.94+4*dx, 0.75), xytext=(1-b0_x, b0_y), **annotate_dict)

ax3.xaxis.set_major_locator(MultipleLocator(0.25))
ax3.yaxis.set_major_locator(MultipleLocator(0.25))
ax3.set_xlabel(r'$t$')
ax3.text(lax, lay, r'(d) $p=8$', **text_dict)
ax3.grid()
plt.setp(ax3.get_yticklabels(), visible=False)

plt.show()

if SERIALIZE:
  fig.savefig("bernstein_extended.pdf", bbox_inches="tight")
\end{minted}

